\documentclass[11pt]{article}
\usepackage{amsmath,amsthm,amssymb}
\usepackage[margin=1in]{geometry}
\usepackage{hyperref}
\usepackage{booktabs}
\usepackage{enumitem}

\newtheorem{theorem}{Theorem}
\newtheorem{lemma}[theorem]{Lemma}
\newtheorem{proposition}[theorem]{Proposition}
\newtheorem{corollary}[theorem]{Corollary}
\newtheorem{conjecture}[theorem]{Conjecture}
\theoremstyle{definition}
\newtheorem{definition}[theorem]{Definition}
\newtheorem{remark}[theorem]{Remark}
\newtheorem{observation}[theorem]{Observation}

\newcommand{\Hq}{H_q}
\newcommand{\Sn}{S_n}
\newcommand{\Vlam}{V^\lambda}
\newcommand{\Om}{\Omega}
\newcommand{\Omtail}{\Om^{(\lambda)}_{\mathrm{tail}}}
\newcommand{\hatl}{\widehat\lambda}
\newcommand{\elh}{\ell(\hatl)}
\newcommand{\Rp}{R'}
\newcommand{\SYT}{\mathrm{SYT}}
\newcommand{\Eplus}[1]{E^+_{#1}}
\newcommand{\Eminus}[1]{E^-_{#1}}
\newcommand{\Vnu}{V^\nu}
\newcommand{\Qsym}{\mathbb{Q}(q)}
\newcommand{\supp}{\mathrm{supp}}

\title{A linear-combination identity for per-subset Pillar~1:\\
       reducing strong operator vanishing to a single tail statement}
\author{Clio}
\date{28 May 2026 (afternoon prove session)}

\begin{document}
\maketitle

\begin{abstract}
For $\lambda$ in the Pillar~1 regime of the multi-row meta-theorem,
the per-subset chain operator $M_T(S_T) := \prod_{k\in\mathrm{tail}} A_k(S_T)$
(with $A_k = P_q^{(i_k)}$ for $k \notin S_T$ and $A_k = P_{-1}^{(i_k)}$
otherwise) is conjecturally contained in $\bigcap_{j=1}^{\tau-|S_T|}\Eminus j|_{\Vlam}$
whenever $|S_T| < \tau$ (\emph{refined per-subset Pillar~1}). We give a
linear-combination identity expressing $M_T(S_T)$ as a signed sum of
\emph{factor-removed tail chains} $\Omtail^{[\widehat U]}$ over $U \subseteq S_T$.
This reduces the per-subset conjecture to a uniform \emph{strong} statement:
$\Omtail^{[\widehat U]}\Vlam \subseteq \Eminus 1|_{\Vlam}$ for every
$U \subseteq \mathrm{tail}$ with $|U| < \tau$. We prove the strong statement
structurally for two distinguished tail positions of $\lambda = (2,2,1,1,1)$
via the $S_{n-1}$-branching method, and verify it computationally for all
$11$ single-factor removals on this shape. The corresponding cases of refined
per-subset Pillar~1 at $|S_T| = 1$ then follow as corollaries of the identity.
Remaining gaps are stated precisely: $9$ tail positions whose structural proof
requires arguments beyond commutation (the braid relation
$S_i S_{i+1} S_i = S_{i+1} S_i S_{i+1}$ fails since
$S_i S_{i+1} S_i - S_{i+1} S_i S_{i+1} = q(T_i - T_{i+1})$);
the $3$ outlier subsets at $|S_T| = \tau$; and the ``$\binom{\ell}{2}-1 \in S$''
sub-case of prefix-kill where refined per-subset Pillar~1 is needed at
strength~$\ge 2$.
\end{abstract}

\section{Setup and statement}\label{sec:setup}

We use the conventions of~\cite{strong}, briefly recalled: $\Hq(\Sn)$ is the
Iwahori--Hecke algebra with $T_i^2 = (q-1)T_i + q$ and braid relations;
$S_i := T_i + 1$ satisfies $S_i^2 = (q+1)S_i$, so $S_i$ has eigenvalues $0, q+1$.
Set $\Eplus i := \ker(T_i - q)|_{\Vlam} = \mathrm{Im}(S_i|_{\Vlam})$ and
$\Eminus i := \ker(S_i|_{\Vlam})$. For $k \ge 1$, write
$\Rp_k := S_{k-1}S_{k-2}\cdots S_1$. The \emph{staircase chain} for $\Sn$ is
$\Om := \Rp_2 \Rp_3 \cdots \Rp_n$, the reduced word for $w_0 \in \Sn$;
total length $N = \binom{n}{2}$.

For $\lambda = (\hatl, 1^q) \vdash n$ with $\elh \ge 2$, every row $\hatl_r
\ge 2$, and $q \ge |\hatl| - 2\elh + 2$, set $\ell := \elh + q$ and
$\tau := q + 2\elh - 1 - |\hatl|$. Split $\Om = R \cdot \Omtail$ with
\[
   R \;:=\; \Rp_2 \Rp_3 \cdots \Rp_\ell,
   \qquad
   \Omtail \;:=\; \Rp_{\ell+1} \Rp_{\ell+2} \cdots \Rp_n.
\]

The \textbf{Pillar~1 meta-theorem}~\cite{pillar1} states
$\Omtail \Vlam \subseteq \bigcap_{j=1}^\tau \Eminus j|_{\Vlam}$.

\paragraph{Per-subset projectors and per-subset chain.}
$P_q^{(i)} := S_i/(q+1)$ projects onto $\Eplus i$;
$P_{-1}^{(i)} := I - P_q^{(i)} = (qI - T_i)/(q+1)$ projects onto $\Eminus i$.
For a subset $S \subseteq \{0, \ldots, N-1\}$ of staircase positions, define
\[
   M(S) \;:=\; \prod_{k=0}^{N-1} A_k(S),
   \qquad
   A_k(S) = \begin{cases} P_{-1}^{(i_k)} & k \in S, \\
                          P_q^{(i_k)} & k \notin S. \end{cases}
\]
Split $S = S_R \sqcup S_T$ with $S_R \subseteq \{0, \ldots, \binom{\ell}{2}-1\}$
(prefix) and $S_T \subseteq \{\binom{\ell}{2}, \ldots, N-1\}$ (tail); set
\[
   M_T(S_T) \;:=\; \prod_{k \in \mathrm{tail}} A_k(S_T),
   \qquad
   M_R(S_R) \;:=\; \prod_{k \in \mathrm{prefix}} A_k(S_R).
\]
Then $M(S) = M_R(S_R) \cdot M_T(S_T)$ and $M_T(\emptyset) = \Omtail/(q+1)^{|\mathrm{tail}|}$.

\paragraph{Factor-removed tail chains.}
For $U \subseteq \mathrm{tail}$, let $\Omtail^{[\widehat U]} \in \Hq$ denote
the product of $S$-factors at the tail positions \emph{not in $U$}:
\[
   \Omtail^{[\widehat U]} \;:=\; \prod_{k \in \mathrm{tail} \setminus U} S_{i_k},
\]
preserving the left-to-right staircase order. Thus
$\Omtail^{[\widehat\emptyset]} = \Omtail$ and
$\Omtail^{[\widehat U]}$ has $|\mathrm{tail}| - |U|$ factors.

\section{The linear-combination identity}\label{sec:identity}

\begin{proposition}[Per-subset $M_T$ as a signed sum of $\Omtail^{[\widehat U]}$]
\label{prop:identity}
For any $S_T \subseteq \mathrm{tail}$,
\begin{equation}\label{eq:identity}
   (q+1)^{|\mathrm{tail}|}\, M_T(S_T) \;=\;
   \sum_{U \subseteq S_T} (-1)^{|S_T|-|U|}\, (q+1)^{|U|}\, \Omtail^{[\widehat U]}.
\end{equation}
\end{proposition}

\begin{proof}
For $k \in S_T$, $A_k = I - P_q^{(i_k)} = I - S_{i_k}/(q+1)$, so we write
\[
   A_k \;=\; \frac{1}{q+1}\bigl((q+1)I - S_{i_k}\bigr),
   \qquad k \in S_T,
\]
while for $k \notin S_T$, $A_k = S_{i_k}/(q+1)$. Therefore
\[
   M_T(S_T) \;=\; \prod_{k} A_k(S_T) \;=\; \frac{1}{(q+1)^{|\mathrm{tail}|}}
   \prod_{k \in \mathrm{tail}}\widetilde A_k,
\]
where $\widetilde A_k = S_{i_k}$ for $k \notin S_T$ and $\widetilde A_k =
(q+1)I - S_{i_k}$ for $k \in S_T$. Expand each $(q+1)I - S_{i_k}$ as a sum
over a binary choice $c_k \in \{0,1\}$: $c_k = 0$ contributes $(q+1)I$ at
position $k$; $c_k = 1$ contributes $-S_{i_k}$. Distributing the product (in
the fixed staircase order) gives
\[
   \prod_{k \in \mathrm{tail}}\widetilde A_k \;=\;
   \sum_{V \subseteq S_T} (-1)^{|V|}\,(q+1)^{|S_T|-|V|}
   \prod_{k}\bigl[I \text{ for } k \in S_T\setminus V,\; S_{i_k}\text{ else}\bigr].
\]
The product on the right is $\Omtail^{[\widehat{(S_T \setminus V)}]}$: we have
$I$ at the $|S_T \setminus V|$ positions in $S_T \setminus V$ (which deletes
those factors) and $S_{i_k}$ at every other position. Substituting
$U := S_T \setminus V$:
\[
   \prod_{k}\widetilde A_k \;=\; \sum_{U \subseteq S_T} (-1)^{|S_T|-|U|}
   (q+1)^{|U|}\, \Omtail^{[\widehat U]},
\]
and dividing by $(q+1)^{|\mathrm{tail}|}$ yields~\eqref{eq:identity}.\qed
\end{proof}

\begin{example}[$|S_T| = 0, 1, 2$]\label{ex:identity-small}
\begin{align*}
   (q+1)^{|\mathrm{tail}|}\, M_T(\emptyset) &= \Omtail,\\
   (q+1)^{|\mathrm{tail}|}\, M_T(\{p\}) &= (q+1)\,\Omtail^{[\widehat{\{p\}}]} - \Omtail,\\
   (q+1)^{|\mathrm{tail}|}\, M_T(\{p,q\}) &=
       (q+1)^2\,\Omtail^{[\widehat{\{p,q\}}]}
       - (q+1)\,\Omtail^{[\widehat{\{p\}}]}
       - (q+1)\,\Omtail^{[\widehat{\{q\}}]}
       + \Omtail.
\end{align*}
\end{example}

\section{Strong Pillar~1 reduction}\label{sec:reduction}

\begin{definition}[Strong Pillar~1]\label{def:strong-pillar1}
Let $\lambda$ be in the Pillar~1 regime. The \emph{Strong Pillar~1
property at level $s$} is the statement
\begin{equation}\label{eq:strong-pillar1}
   \Omtail^{[\widehat U]} \Vlam \;\subseteq\;
   \Eminus 1|_{\Vlam}
   \qquad \forall\, U \subseteq \mathrm{tail} \text{ with } |U| \le s.
\end{equation}
\end{definition}

For $s = 0$, this is Pillar~1 (whose conclusion is in fact stronger:
$\subseteq \bigcap_{j=1}^\tau \Eminus j$). For $s \ge 1$ it is a
nontrivial extension.

\begin{theorem}[Per-subset Pillar~1 from Strong Pillar~1]\label{thm:per-subset-from-strong}
If the Strong Pillar~1 property at level $s$ holds (Definition~\ref{def:strong-pillar1}),
then for every $S_T \subseteq \mathrm{tail}$ with $|S_T| \le s$,
\[
   M_T(S_T) \Vlam \;\subseteq\; \Eminus 1|_{\Vlam}.
\]
\end{theorem}

\begin{proof}
By Proposition~\ref{prop:identity},
\[
   (q+1)^{|\mathrm{tail}|}\, M_T(S_T) \Vlam
   \;\subseteq\; \sum_{U \subseteq S_T} \Omtail^{[\widehat U]} \Vlam
   \;\subseteq\; \Eminus 1
\]
where the last inclusion uses the Strong Pillar~1 hypothesis for every
$U \subseteq S_T$, each of size $\le s$.\qed
\end{proof}

\begin{remark}[Refined per-subset Pillar~1]\label{rem:refined}
The stronger refinement
$M_T(S_T) \Vlam \subseteq \bigcap_{j=1}^{\tau-|S_T|}\Eminus j$
(Conjecture~3 of~\cite{strong}) likewise reduces to a \emph{refined Strong Pillar~1}:
$\Omtail^{[\widehat U]} \Vlam \subseteq \bigcap_{j=1}^{\tau-|U|}\Eminus j$
for every $U \subseteq \mathrm{tail}$ with $|U| < \tau$, by the same proof.
At $\lambda = (2,2,1,1,1)$ ($\tau = 2$) this gives, at $|U| = 1$, the bound
$\Omtail^{[\widehat U]} \Vlam \subseteq \Eminus 1$, matching
Definition~\ref{def:strong-pillar1}.
\end{remark}

\section{Strong Pillar~1 at $\lambda = (2,2,1,1,1)$: two structural proofs}
\label{sec:structural}

We focus on $\lambda = (2,2,1,1,1)$ where $\hatl = (2,2)$, $\elh = 2$, $q = 3$,
$\ell = 5$, $\tau = 2$, $n = 7$, $N = 21$. The tail consists of positions
$10$ through $20$ corresponding to staircase letters
$(5, 4, 3, 2, 1; 6, 5, 4, 3, 2, 1)$: blocks $\Rp_6 = S_5 S_4 S_3 S_2 S_1$
(positions $10$--$14$) and $\Rp_7 = S_6 S_5 S_4 S_3 S_2 S_1$ (positions
$15$--$20$).

The $S_{n-1} = S_6$ branching of $\Vlam$ is
\[
   \Vlam \downarrow S_6 \;\cong\; V^{(2,1,1,1,1)} \oplus V^{(2,2,1,1)},
\]
where we denote the summands $W_1, W_2$ respectively.

\begin{lemma}[Structural Strong Pillar~1 at position $10$]\label{lem:pos10}
$\Omtail^{[\widehat{\{10\}}]} \Vlam \subseteq \Eminus 1|_{\Vlam}$.
\end{lemma}

\begin{proof}
Position $10$ corresponds to the leftmost $S_5$ of $\Rp_6$. Removing it leaves
$\Rp_6^{[\widehat{S_5}]} = S_4 S_3 S_2 S_1 = \Rp_5$, so
\[
   \Omtail^{[\widehat{\{10\}}]} \;=\; \Rp_5 \cdot \Rp_7
   \;=\; \Rp_5 \cdot (S_6\, \Rp_6),
\]
using $\Rp_7 = S_6\, \Rp_6$.

\emph{Step 1: Commutation.} The letters of $\Rp_5 = S_4 S_3 S_2 S_1$ are all
$\le 4$, so $|6 - i| \ge 2$ for each $i$ and $S_6$ commutes with every factor
of $\Rp_5$. Therefore $\Rp_5 \cdot S_6 = S_6 \cdot \Rp_5$, giving
\[
   \Omtail^{[\widehat{\{10\}}]} \;=\; S_6 \cdot (\Rp_5 \Rp_6).
\]

\emph{Step 2: $\Rp_5 \Rp_6$ acts via $S_6$-branching.} Since
$\Rp_5 \Rp_6 \in \Hq(S_6)$ (no $T_6$), the operator preserves the $W_1 \oplus W_2$
decomposition.

\smallskip\emph{On $W_2 = V^{(2,2,1,1)}$:} the shape $\nu = (2,2,1,1)$ has
$\widehat\nu = (2,2)$, $\elh = 2$, $q'' = 2$, satisfying the Pillar~1 hypotheses
($q'' \ge |\widehat\nu| - 2\elh + 2 = 2$, $\tau(\widehat\nu) = 1$). The chain
$\Rp_5\Rp_6 \in \Hq(S_6)$ is exactly $\Omtail^{(2,2,1,1)}$ (blocks at indices
$\ell(\nu)+1 = 5$ and $\ell(\nu)+2 = 6$), so Pillar~1 applies:
\[
   \Rp_5 \Rp_6\, W_2 \;\subseteq\; \Eminus 1|_{W_2}.
\]

\smallskip\emph{On $W_1 = V^{(2,1,1,1,1)}$:} this is a hook with
$k_h = 2$, $m_h = 4$, satisfying $m_h \ge k_h$. The hook column-$1$
theorem~\cite[Thm.~6]{pillar1} gives $\Rp_6\, W_1 \subseteq \Eminus 1|_{W_1}$.
Then $\Rp_5 = S_4 S_3 S_2 S_1$ has rightmost factor $S_1$, which is identically
zero on $\Eminus 1 = \ker(S_1)$:
\[
   \Rp_5 \cdot \Rp_6 W_1 \;=\; S_4 S_3 S_2 \cdot (S_1 \cdot \Eminus 1|_{W_1})
   \;=\; 0.
\]

\smallskip Combining,
$\Rp_5\Rp_6\,\Vlam = \Rp_5\Rp_6(W_1) \oplus \Rp_5\Rp_6(W_2)
\subseteq 0 \oplus \Eminus 1|_{W_2} \subseteq \Eminus 1|_{\Vlam}$.

\emph{Step 3: $S_6$ preserves $\Eminus 1$.} Since $|6 - 1| = 5 \ge 2$, $S_6$
commutes with $T_1$ and therefore preserves $\Eminus 1$. Hence
\[
   \Omtail^{[\widehat{\{10\}}]}\Vlam \;=\; S_6 \cdot \Rp_5 \Rp_6\, \Vlam
   \;\subseteq\; S_6\, \Eminus 1 \;\subseteq\; \Eminus 1.\qed
\]
\end{proof}

\begin{lemma}[Structural Strong Pillar~1 at position $15$]\label{lem:pos15}
$\Omtail^{[\widehat{\{15\}}]} \Vlam \subseteq \Eminus 1|_{\Vlam}$.
\end{lemma}

\begin{proof}
Position $15$ corresponds to the leftmost $S_6$ of $\Rp_7$. Removing it leaves
$\Rp_7^{[\widehat{S_6}]} = S_5 S_4 S_3 S_2 S_1 = \Rp_6$, so
\[
   \Omtail^{[\widehat{\{15\}}]} \;=\; \Rp_6 \cdot \Rp_6 \;=\; (\Rp_6)^2.
\]

\emph{Step 1: Algebraic identity.} Using $\Rp_6 = S_5 \cdot \Rp_5$
($\Rp_6 = S_5 S_4 S_3 S_2 S_1$), associativity gives
\[
   (\Rp_6)^2 \;=\; (S_5 \Rp_5)\cdot \Rp_6 \;=\; S_5 \cdot \Rp_5 \Rp_6.
\]

\emph{Step 2: $\Rp_5\Rp_6 \in \Hq(S_6)$ and the same branching argument as
Lemma~\ref{lem:pos10}} gives $\Rp_5\Rp_6\,\Vlam \subseteq \Eminus 1|_{\Vlam}$.

\emph{Step 3: $S_5$ preserves $\Eminus 1$.} Since $|5 - 1| = 4 \ge 2$, $S_5$
commutes with $T_1$ and therefore preserves $\Eminus 1$. Hence
\[
   \Omtail^{[\widehat{\{15\}}]}\Vlam \;=\; S_5 \cdot \Rp_5\Rp_6\,\Vlam
   \;\subseteq\; S_5\, \Eminus 1 \;\subseteq\; \Eminus 1.\qed
\]
\end{proof}

\begin{corollary}[Per-subset Pillar~1 at $|S_T| = 1$ for positions $10, 15$]
\label{cor:pos-10-15-per-subset}
For $S_T \in \{\{10\}, \{15\}\}$,
\[
   M_T(S_T)\,\Vlam \;\subseteq\; \Eminus 1|_{\Vlam}.
\]
\end{corollary}
\begin{proof}
By Theorem~\ref{thm:per-subset-from-strong} and Lemmas~\ref{lem:pos10}--\ref{lem:pos15}.
\qed
\end{proof}

\section{Computational verification of Strong Pillar~1}\label{sec:verify}

We verify the full Strong Pillar~1 property at $\lambda = (2,2,1,1,1)$ for
all single-factor removals and identify the residual cases at pair removals.
Computations use exact rationals at $q = 5/7$, cross-checked at $q = 2/3$.

\begin{table}[h]
\centering\renewcommand{\arraystretch}{1.2}
\caption{$\Omtail^{[\widehat{\{p\}}]}\Vlam$ for each tail position $p$
on $\lambda = (2,2,1,1,1)$. Rank, $\Eminus 1$ and $\Eminus 2$ containment
verified by exact computation.}
\begin{tabular}{ccccccc}
\toprule
position $p$ & letter $i_p$ & block & rank of $\Omtail^{[\widehat{\{p\}}]}|_{\Vlam}$ & $\subseteq \Eminus 1$ & $\subseteq \Eminus 2$ & structural proof? \\
\midrule
$10$ & $5$ & $\Rp_6$ (leftmost) & $1$ & yes & yes & Lemma~\ref{lem:pos10} \\
$11$ & $4$ & $\Rp_6$ & $1$ & yes & no & empirical only \\
$12$ & $3$ & $\Rp_6$ & $0$ & yes & yes & empirical only \\
$13$ & $2$ & $\Rp_6$ & $0$ & yes & yes & empirical only \\
$14$ & $1$ & $\Rp_6$ (rightmost) & $1$ & yes & yes & empirical only \\
$15$ & $6$ & $\Rp_7$ (leftmost) & $1$ & yes & no & Lemma~\ref{lem:pos15} \\
$16$ & $5$ & $\Rp_7$ & $0$ & yes & yes & empirical only \\
$17$ & $4$ & $\Rp_7$ & $0$ & yes & yes & empirical only \\
$18$ & $3$ & $\Rp_7$ & $1$ & yes & yes & empirical only \\
$19$ & $2$ & $\Rp_7$ & $1$ & yes & yes & empirical only \\
$20$ & $1$ & $\Rp_7$ (rightmost) & $1$ & yes & yes & empirical only \\
\bottomrule
\end{tabular}
\end{table}

All $11$ single-removals lie in $\Eminus 1$; $9$ of them also lie in
$\Eminus 2$ (the sharper bound). Combining with the linear-combination
identity (Proposition~\ref{prop:identity}), we obtain
$M_T(\{p\})\Vlam \subseteq \Eminus 1$ for every singleton $S_T$.

For pair removals $|U| = 2$, $52$ of the $\binom{11}{2} = 55$ pairs satisfy
$\Omtail^{[\widehat U]}\Vlam \subseteq \Eminus 1$; the remaining $3$
\emph{outliers} are $U \in \{\{11,12\}, \{12,15\}, \{15,16\}\}$, matching
exactly the three $|S_T| = \tau$ outlier subsets identified in~\cite{strong}.
At $|S_T| = 2 = \tau$, refined per-subset Pillar~1 makes no claim, so the
outliers are outside the scope of Theorem~\ref{thm:per-subset-from-strong}
(level $s = 1$).

\section{Why braid does not help (and what does)}\label{sec:braid}

A natural temptation is to use the Coxeter braid relation $S_i S_{i+1} S_i
= S_{i+1} S_i S_{i+1}$ to manipulate tail chains, e.g., to extend the
Lemma~\ref{lem:pos10} commutation argument to position $11$. This fails for
$S_i = T_i + 1$:

\begin{lemma}[$S$-braid defect]\label{lem:braid-defect}
In $\Hq$, $S_i S_{i+1} S_i - S_{i+1} S_i S_{i+1} = q\,(T_i - T_{i+1})$.
\end{lemma}

\begin{proof}
Expanding $(T_i+1)(T_{i+1}+1)(T_i+1)$ and using $T_i^2 = (q-1)T_i + q$:
\[
   S_i S_{i+1} S_i \;=\; T_i T_{i+1} T_i + T_i T_{i+1} + T_{i+1} T_i + (q+1) T_i + T_{i+1} + (q+1).
\]
Symmetrically,
$S_{i+1} S_i S_{i+1} = T_{i+1} T_i T_{i+1} + T_{i+1} T_i + T_i T_{i+1} + (q+1) T_{i+1} + T_i + (q+1)$.
Subtracting (using $T_i T_{i+1} T_i = T_{i+1} T_i T_{i+1}$) yields
$(q+1)(T_i - T_{i+1}) - (T_i - T_{i+1}) = q(T_i - T_{i+1})$.\qed
\end{proof}

\begin{remark}[Why positions $10$ and $15$ are special]\label{rem:special-positions}
The proofs of Lemmas~\ref{lem:pos10}--\ref{lem:pos15} use only:
(a) the commutation $S_i S_j = S_j S_i$ for $|i - j| \ge 2$ (Coxeter commutation);
and (b) the trivial associativity identity $S_5\, \Rp_5 = \Rp_6$.
No braid relation is invoked. The two privileged positions are exactly the
\emph{leftmost} factors of $\Rp_6$ and $\Rp_7$ (positions $10$ and $15$,
letters $5$ and $6$), because removing the leftmost $S_5$ of $\Rp_6$ leaves
exactly $\Rp_5$, and removing the leftmost $S_6$ of $\Rp_7$ leaves
exactly $\Rp_6$. Both produce a standard $\Omtail$-like chain (for a smaller
shape or by squaring) that succumbs to Pillar~1 plus $S_1$-kill on the
$W_1$-summand. The remaining $9$ positions yield modified chains that are not
of this standard form; structural proofs would require either an enhanced
form of Pillar~1 or genuine use of the $S$-braid defect from
Lemma~\ref{lem:braid-defect} (with the correction term integrated into the
analysis).
\end{remark}

\section{Status and gaps}\label{sec:status}

\paragraph{Established.}
\begin{itemize}[topsep=0.2em,itemsep=0.1em]
\item Proposition~\ref{prop:identity}: linear-combination identity expressing
   $M_T(S_T)$ as a signed sum of factor-removed tail chains, with coefficients
   $\pm(q+1)^{|U|}$.
\item Theorem~\ref{thm:per-subset-from-strong}: per-subset Pillar~1 at level
   $s$ reduces to Strong Pillar~1 at level $s$, a uniform statement depending
   only on $|U|$ and not on the choice of $P_q$ vs.\ $P_{-1}$ at each position.
\item Lemmas~\ref{lem:pos10} and~\ref{lem:pos15}: Strong Pillar~1 at the
   distinguished positions $10$ and $15$ at $\lambda = (2,2,1,1,1)$, by an
   $S_{n-1}$-branching argument using only commutation and Pillar~1 on smaller
   shapes.
\item Corollary~\ref{cor:pos-10-15-per-subset}: per-subset Pillar~1 at
   $|S_T| = 1$ for the two singleton subsets at positions $10$ and $15$.
\item Lemma~\ref{lem:braid-defect}: explicit $q$-correction in the
   $S$-braid identity.
\end{itemize}

\paragraph{Open gaps (precisely stated).}
\begin{itemize}[topsep=0.2em,itemsep=0.1em]
\item Strong Pillar~1 at positions $11, 12, 13, 14, 16, 17, 18, 19, 20$ (nine
   positions) for $\lambda = (2,2,1,1,1)$: verified computationally; structural
   proof open. Of these, positions $12, 13, 16, 17$ have $\Omtail^{[\widehat
   {\{p\}}]} = 0$ as an operator on $\Vlam$ (a strictly stronger statement
   than $\subseteq \Eminus 1$), suggesting they have their own clean argument
   we have not yet identified.
\item Strong Pillar~1 at level $s = 2$ for $52$ of the $55$ pair-removals.
   The $3$ outliers correspond to $|S_T| = \tau$ and lie outside refined
   per-subset Pillar~1; they require a separate inductive argument on
   $\Om_{S_\ell}$-kernels in branching summands (the residual gap
   from~\cite{strong}).
\item The prefix-kill cases where $\binom{\ell}{2}-1 \in S$ (position $9$ in
   $S$): refined per-subset Pillar~1 at level $\ge 2$ would close these.
\item Strong Pillar~1 at general $\lambda = (2,2,1^m)$ for $m \ge 4$:
   computational verification beyond $m = 3$ is open, as is the inductive
   structure of the argument across $m$.
\end{itemize}

\paragraph{Methodological note.}
The two structural Lemmas use the $S_{n-1}$-branching method (Lemma~\ref{lem:pos10}
Step~2): an operator that lies in $\Hq(S_{n-1})$ acts on $\Vlam \downarrow
S_{n-1}$ summand-by-summand. On each summand $V^\nu$, smaller-shape Pillar~1
or the hook column-$1$ theorem applies. The crucial input is that the
\emph{leftmost} factor of the tail (the $S_5$ at position $10$, or the
$S_6$ at position $15$) is the one whose removal yields a chain in $\Hq(S_{n-1})$.
For positions deeper in $\Rp_6$ or $\Rp_7$, the resulting chain contains
$S_5$ or $S_6$ entangled with the lower letters, and the $\Hq(S_{n-1})$
factorisation requires a braid-like manipulation that introduces the
$q(T_i - T_{i+1})$ correction of Lemma~\ref{lem:braid-defect}. We have not
found a clean way to integrate this correction into the branching argument
in this session.

\begin{thebibliography}{9}
\bibitem{strong}
   Clio, \emph{Strong operator vanishing of the staircase chain on $V^{(2,2,1^m)}$
   (the $|S|=0$ case of strong per-subset vanishing)}, 28 May 2026.
   \texttt{\textasciitilde/projects/proofs/2026-05-28-strong-per-subset-operator-vanishing.tex}.

\bibitem{pillar1}
   Clio, \emph{Extended sign-kill for general multi-row $\hatl$: the
   meta-theorem at arbitrary $\elh \ge 2$}, 15 May 2026.
   \texttt{\textasciitilde/projects/proofs/2026-05-15-multi-row-sign-kill-meta.tex}.

\bibitem{sharpness}
   Clio, \emph{Sharpness of the multi-row containment at $\tau+1$:
   $B^{(\lambda)} \cap E^-_{\tau+1}|_{\Vlam} = 0$}, 17 May 2026.
   \texttt{\textasciitilde/projects/proofs/2026-05-17-sharpness-tau-plus-1.tex}.
\end{thebibliography}

\end{document}
